\section{Digital Forensics and Generative AI}

The integration of LLMs into digital forensics offers new ways to speed up evidence analysis.\\
However, because LLMs operate on probabilistic principles, their application introduces unique risks concerning evidence validation, repeatability, and legal admissibility.

\subsection{LLM Workflows}

LLMs are advanced AI systems that generate human-like text outputs.\\
They utilize trillions of parameters to compute linguistic relationships between tokens regardless of their distance within a dataset.

\subsubsection*{Enhancing Efficiency across Investigative Phases}
When applied to data, generative models streamline the parsing of unstructured textual or visual sources:
\begin{itemize}
    \item \textbf{Early Log and Data Analysis:} LLMs parse large volumes of unstructured or semi-structured system logs to flag anomalies, detect unauthorized access attempts, identify potential security breaches, and map early Indicators of Compromise (IoC).
    \item \textbf{Textual and Visual Evidence Sorting:} Advanced foundation models (such as GPT-4) excel at analyzing raw textual content, interpreting written evidence, and drafting preliminary forensic reports.
    \item \textbf{Multimodal Contextual Analysis:} Specialized visual models combine Natural Language Processing (NLP) with computer vision to automatically detect and categorize objects or faces inside images
\end{itemize}

\subsection{Operational Risks}

Despite their operational utility, generative frameworks pose distinct structural challenges when evaluated against strict forensic standards.

\subsubsection*{1. Reliability and Adversarial Vulnerabilities}
The reliability of AI output is directly correlated with anti-forensic countermeasures.
\begin{itemize}
    \item \textbf{Adversarial Robustness Flaws:} Bad actors can evade detection by altering file extensions, embedding steganographic data, or using adversarial perturbations specifically designed to trick AI classifiers.
    \item \textbf{Visual Patch Vulnerabilities:} Tiny pixel alterations, invisible to a human investigator, can cause an AI to misclassify an illicit item.
\end{itemize}

\subsubsection*{2. Repeatability vs. Probabilistic Incoherence}
Digital forensics requires strict determinism, meaning that executing the exact same analytical process over identical evidence must always yield an identical result.
\begin{itemize}
    \item \textbf{Probabilistic Problem:} Because LLMs and agentic AI frameworks run on probabilistic predictions, entering the exact same prompt twice through a pipeline (such as an AutoGen setup) can yield two different execution paths, text summaries, or script outputs.
    \item \textbf{Solutions:} Current research focuses on integrating deterministic sandboxing and verification tools. Frameworks like \textit{AutoDFBench} have been developed to rigorously test and validate AI-generated forensic scripts against established NIST standards to ensure compliance.
\end{itemize}

\subsubsection*{3. Recoverability}
The process of file reconstructing corrupt data introduces a shift from physical byte validation to predictive generation.
\begin{itemize}
    \item \textbf{Recovery Target}
    \begin{itemize}
        \item \textbf{Traditional Data Carving:} Searches for rigid binary signatures, such as file headers and footers (e.g., \texttt{\textbackslash xFF\textbackslash xD8\textbackslash xFF} for JPEGs).
        \item \textbf{AI-Assisted Semantic Carving:} Evaluates the linguistic, semantic, and contextual flow between unallocated text fragments.
    \end{itemize}

    \item \textbf{Execution Method}
    \begin{itemize}
        \item \textbf{Traditional Data Carving:} Physically maps and stitches continuous raw disk sectors back together.
        \item \textbf{AI-Assisted Semantic Carving:} Predicts how missing or damaged data strings should logically connect based on learned linguistic patterns.
    \end{itemize}

    \item \textbf{Legal Admissibility}
    \begin{itemize}
        \item \textbf{Traditional Data Carving:} Data either matches structural rules or fails.
        \item \textbf{AI-Assisted Semantic Carving:} ùThe resulting file may appear correct but cuold be wrong, making it easy to challenge in court.
    \end{itemize}
\end{itemize}